\section{From the Model to Real Chess: Details}\label{S-real}
We build the tree under either passed-pawn test, the \emph{textbook tree} or the \emph{e.p.-aware tree}.

\begin{theorem}[Candidate in real chess; \lab{C1}]
For the candidate \pos{W a2 c2 / B b7 d7} beside T (Fig.~\ref{fig:star2_textbook}):
\begin{enumerate}
\item[(a)] In chess the side to move wins; with black to move only 1\ldots~b5 wins; with white to
  move 1.~c4 wins, 1.~a3, 1.~c3 and all king moves lose, and 1.~a4 is unresolved \lab{C2}.
\item[(b)] In the e.p.-aware tree, for every draw treatment, $O$ has left stop $\ge1$ and right stop $\le-1$, so
  it is not $0$ or $\nim{2}$.
\item[(c)] With draws as wins for black, $O=\game{2\pm1}{-2\pm1}$ in both trees.
\end{enumerate}
\end{theorem}

Real chess thus agrees with $\nim{2}$ on the outcome and on every winning first move resolved so far
(1.~a4 is still open), but not on the value: in the e.p.-aware tree, which matches chess, $O$ is hot under every draw treatment, so
the candidate's $\nim{2}$ is a feature of the textbook test, not of chess. In the textbook tree we rule out
$\nim{2}$ only when draws count as wins for black, as in (c), or for white \lab{C1}; with draws worth $0$ we do
not decide it.

A second candidate, \pos{W a3 b3 c2 / B b5 b7 c6} ($10$ men with T), is $\nim{2}$ under the opponent
rule with en passant on (Section~\ref{sec:star2}); its oracle value is positive under every draw treatment
\lab{C1}. The other non-default $\nim{2}$ positions are open.

\subsection{Towards a Real-Chess \texorpdfstring{$1/8$}{1/8}}

The $1/8$ position of Section~\ref{sec:targets} (Fig.~\ref{fig:eighth}(a)) has $10$ men. Its e.p.-aware tree
has $302$ end-position sides with at most seven men, and the tablebase resolves them all. In $290$ of them
only one side has a passed pawn: $283$ agree with the model, $5$ disagree and $2$ are draws; the other $12$ are
races. A search for the larger end positions that decide the value finds five with eight men and three with
nine. Four of the eight-man ones are resolved: seven of their sides agree with the model, and black to move in
\pos{W a3 a5 / B b4 b5} is a black win against it. A second search allows for that result and leaves one
eight-man position and the three of nine men to decide the value. Our king-and-pawn solver (queen promotions
only) shows that all four agree with the model on both sides; the hardest, black to move in
\pos{W b2 b5 / B a5 b4}, needed two promoted queens in the search, $6.1\times10^{10}$ states in $15.6$~h.

\begin{proposition}[A tablebase-backed oracle value of $1/8$; \lab{C1}]
For the $1/8$ position, with queen promotions only in our solver, $O=1/8$, whatever the open end positions do.
\end{proposition}

This $1/8$ assumes the frozen kings of the tree (Section~\ref{sec:kings}); whether the full-chess position
has value $1/8$ is open. The same method applies to $\upsec\star$ and $\tinyg{1}$: their lower bounds already equal the
target values, with $10$ and $12$ end positions of $9$ or $10$ men still open \lab{C1}. Positions A and B have $11$ men beside
T, and the five-file $\nim{2}$ \pos{W a3 b3 c2 e2 / B a6 c5 d5 e3}, which sits legally
beside T, has $12$; no file between its pawns and T is free of pawns, so Theorem~\ref{thm:validity} does not apply. Their oracle values are open: for A and B a
search finds the deciding end positions; the five and seven of them with $8$ men agree with the model
\lab{C2}, but those with $9$ to $11$ men are unsolved, so their $\nim{2}$ is a model value.
